% Verify every entry (authors, volume, pages, DOI) before submission.
\begin{thebibliography}{99}\small
\bibitem{iea2023} International Energy Agency, Energy Efficiency 2023, IEA, Paris, 2023 (section: Can efficient cooling help manage fast rising electricity demand in India and achieve thermal comfort for all?).
\bibitem{drgona2020} J.~Drgo\v{n}a, J.~Arroyo, I.~Cupeiro Figueroa, D.~Blum, K.~Arendt, D.~Kim, E.~Perarnau Oll\'e, J.~Oravec, M.~Wetter, D.L.~Vrabie, L.~Helsen, All you need to know about model predictive control for buildings, Annual Reviews in Control 50 (2020) 190--232.
\bibitem{yao2021} Y.~Yao, D.K.~Shekhar, State of the art review on model predictive control (MPC) in heating ventilation and air-conditioning (HVAC) field, Building and Environment 200 (2021) 107952.
\bibitem{yu2021} L.~Yu, S.~Qin, M.~Zhang, C.~Shen, T.~Jiang, X.~Guan, A review of deep reinforcement learning for smart building energy management, IEEE Internet of Things Journal 8 (2021) 12046--12063.
\bibitem{sierla2022} S.~Sierla, H.~Ihasalo, V.~Vyatkin, A review of reinforcement learning applications to control of heating, ventilation and air conditioning systems, Energies 15 (2022) 3526.
\bibitem{arroyo2022} J.~Arroyo, C.~Manna, F.~Spiessens, L.~Helsen, Reinforced model predictive control (RL-MPC) for building energy management, Applied Energy 309 (2022) 118346.
\bibitem{mnih2015} V.~Mnih et al., Human-level control through deep reinforcement learning, Nature 518 (2015) 529--533.
\bibitem{vanhasselt2016} H.~van Hasselt, A.~Guez, D.~Silver, Deep reinforcement learning with double Q-learning, Proc. AAAI Conference on Artificial Intelligence 30 (2016).
\bibitem{alshiekh2018} M.~Alshiekh, R.~Bloem, R.~Ehlers, B.~K\"onighofer, S.~Niekum, U.~Topcu, Safe reinforcement learning via shielding, Proc. AAAI Conference on Artificial Intelligence 32 (2018) 2669--2678. doi:10.1609/aaai.v32i1.11797.
\bibitem{skogestad2003} S.~Skogestad, Simple analytic rules for model reduction and PID controller tuning, Journal of Process Control 13 (2003) 291--309.
\bibitem{beck2009} A.~Beck, M.~Teboulle, A fast iterative shrinkage-thresholding algorithm for linear inverse problems, SIAM Journal on Imaging Sciences 2 (2009) 183--202.
\bibitem{wang2024} X.~Wang, B.~Dong, Long-term experimental evaluation and comparison of advanced controls for HVAC systems, Applied Energy 371 (2024) 123706.
\bibitem{biemann2021} M.~Biemann, F.~Scheller, X.~Liu, L.~Huang, Experimental evaluation of model-free reinforcement learning algorithms for continuous HVAC control, Applied Energy 298 (2021) 117164.
\bibitem{manu2016} S.~Manu, Y.~Shukla, R.~Rawal, L.E.~Thomas, R.~de~Dear, Field studies of thermal comfort across multiple climate zones for the subcontinent: India Model for Adaptive Comfort (IMAC), Building and Environment 98 (2016) 55--70.
\bibitem{nbc2005} Bureau of Indian Standards, National Building Code of India 2005, Part 8 Building Services, Section 3: Air conditioning, heating and mechanical ventilation, Cl.~4.4.3 (Inside design conditions), Table~2 (summer design range for air-conditioned occupancies 23--26\,$^\circ$C DB).
\bibitem{ens2018} Bureau of Energy Efficiency, Eco-Niwas Samhita 2018 (Energy Conservation Building Code for Residential Buildings), Part I: Building Envelope.
\bibitem{ashrae622} ASHRAE, ANSI/ASHRAE Standard 62.2, Ventilation and Acceptable Indoor Air Quality in Residential Buildings (whole-dwelling rate $Q_{tot}=0.03A_{floor}+7.5(N_{br}+1)$\,cfm, unchanged since the 2013 edition).
\bibitem{belias2022} E.~Belias, D.~Licina, Outdoor PM2.5 air filtration: optimising indoor air quality and energy, Buildings \& Cities 3 (2022) 186--203. doi:10.5334/bc.153.
\bibitem{who2021} World Health Organization, WHO global air quality guidelines: particulate matter (PM2.5 and PM10), ozone, nitrogen dioxide, sulfur dioxide and carbon monoxide, Geneva, 2021.
\bibitem{cpcb2009} Central Pollution Control Board, National Ambient Air Quality Standards, Notification, 18 November 2009.
\bibitem{merc} Maharashtra Electricity Regulatory Commission, Multi-Year Tariff Order for MSEDCL, Case No.~217 of 2024 (28 March 2025), as modified by the review order of 25 June 2025 (residential ToD rebate of INR 0.80/kWh during solar hours 09:00--17:00 in FY~2025--26, rising to INR 1.00/kWh by FY~2029--30; night rebate removed).
\bibitem{cea} Central Electricity Authority, CO$_2$ Baseline Database for the Indian Power Sector, Version 20.0, User Guide, December 2024 (weighted average emission factor of the Indian grid, FY2023-24: 0.727 tCO$_2$/MWh).
\bibitem{ishrae} ISHRAE, Indian Weather Data for Energy Simulation (ISHRAE/EnergyPlus weather files), Pune 430630.
\end{thebibliography}
